\section{张量代数的函子性}

记号约定:$A$是交换环.

\begin{proposition}{张量代数的自然性}{naturality-of-tensor-algebra}
	设$M,N$是$A$-模,$u\in\Hom_{A}\left(M,N\right)$,则存在唯一的$u^{\prime}\in\Hom_{A-\mathsf{Alg}}\left(\mathsf{T}\left(M\right),\mathsf{T}\left(N\right)\right)$使下图交换.
	\begin{center}
		\begin{tikzcd}
			M & N \\
			{\mathsf{T}\left(M\right)} & {\mathsf{T}\left(N\right)}
			\arrow["u", from=1-1, to=1-2]
			\arrow["{\phi_{M}}"', from=1-1, to=2-1]
			\arrow["{\phi_{N}}", from=1-2, to=2-2]
			\arrow["{\exists!u^{\prime}}"', dashed, from=2-1, to=2-2]
		\end{tikzcd}
	\end{center}
\end{proposition}

\begin{definition}{映射的典范扩张}{}
	沿用命题(\ref{prop:naturality-of-tensor-algebra})的设定.我们把$u^{\prime}$记作$\mathsf{T}\left(u\right)$,称之为$u$到$T\left(M\right)$的典范扩张.
\end{definition}

\begin{proposition}{张量代数的函子性}{}
	设$M,N$是$A$-模,$u\in\Hom_{A}\left(M,N\right),v\in\Hom_{A}\left(N,P\right)$,则下列图表交换.
	\begin{center}
		\begin{tikzcd}
			M & N & P \\
			{\mathsf{T}(M)} & {\mathsf{T}(N)} & {\mathsf{T}(N)}
			\arrow["u", from=1-1, to=1-2]
			\arrow["{\mathsf{T}(v\circ u)}", curve={height=-18pt}, from=1-1, to=1-3]
			\arrow["{\phi_{M}}"', from=1-1, to=2-1]
			\arrow["v", from=1-2, to=1-3]
			\arrow["{\phi_{N}}", from=1-2, to=2-2]
			\arrow["{\phi_{P}}", from=1-3, to=2-3]
			\arrow["{\mathsf{T}(u)}"', from=2-1, to=2-2]
			\arrow["{\mathsf{T}(v)\circ\mathsf{T}(u)}"', curve={height=24pt}, from=2-1, to=2-3]
			\arrow["{\mathsf{T}(v)}"', from=2-2, to=2-3]
		\end{tikzcd}
	\end{center}
\end{proposition}

\begin{remark}{}{}
	设$n\in\N$.记$\mathsf{T}^{n}\left(u\right):\mathsf{T}^{n}\left(M\right)\to\mathsf{T}^{n}\left(N\right),x\mapsto T\left(u\right)\left(x\right)$,那么$\mathsf{T}^{n}\left(u\right)=u^{\otimes n}$.我们称$\mathsf{T}^{n}\left(u\right)$是$u$的$n$次张量幂.
\end{remark}

\begin{proposition}{$n$次张量幂的性质}{}
	设$M,N$是$A$-模,$u\in\Hom_{A}\left(M,N\right)$.
	\begin{enumerate}
		\item 若$u$是满射,则$\mathsf{T}\left(u\right)$是满射.
		\item $\ker\left(\mathsf{T}\left(u\right)\right)$是$\ker\left(u\right)$在$\mathsf{T}\left(M\right)$中生成的双边理想.
		\item $\mathsf{T}^{0}\left(u\right)$是双射,对每个$n>0$,$\mathsf{T}^{n}\left(u\right)$是满射.
	\end{enumerate}
\end{proposition}

\begin{remark}{}{}
	如果$u$是单射,那么$\mathsf{T}\left(u\right)$一般不是单射.
\end{remark}

\begin{proposition}{}{}
	设$M$是$A$-模,$N,P$是$M$的子模,$N+P$是$M$的直因子,$N\cap P$是$N$和$P$各自的直因子.考虑典范单射
	\begin{align*}
		i:N\to M,j:P\to M,k:N\cap P\to M.
	\end{align*}
	\begin{enumerate}
		\item $\mathsf{T}\left(i\right),\mathsf{T}\left(j\right),\mathsf{T}\left(k\right)$都是单射.
		\item $\im\left(\mathsf{T}\left(k\right)\right)=\im\left(\mathsf{T}\left(i\right)\right)\cap\im\left(\mathsf{T}\left(j\right)\right)$.
	\end{enumerate}
\end{proposition}

\begin{remark}{}{}
	如果$A$是域,那么``$N\cap P$是$N$和$P$各自的直因子''自动成立.此外,如果$N\subsetneq P$,则$\forall n\geq 1\left(\mathsf{T}^{n}\left(N\right)\neq\mathsf{T}^{n}\left(P\right)\right)$.
\end{remark}























